%%
%% Chapter 5: Conclusion
%% Independent file, included by ../main.tex
%%

\section{Conclusion}
\label{sec:conclusion}

We presented a probabilistic redesign of Gaussian-splatting-based
sparse-view CBCT reconstruction built on a latent neural process, with
particular emphasis on the P5 GT-conditioned CVAE variant. P5 borrows
ground-truth information during training through a lightweight 3D CNN
posterior encoder, yet discards both the GT encoder and the dense views
at deployment, sampling exclusively from a sparse-view prior. On the
LUNA16 benchmark with six projections, P5 reaches $29.24$\,dB PSNR /
$0.863$\,SSIM (best-eval), within $0.3$\,dB of the deterministic
reference, while providing per-point predictive variance at no
additional inference cost.

The central methodological contribution of this paper is diagnostic
rather than architectural. Through systematic instrumentation -- KL
decomposition, cross-case latent statistics, prior/posterior
reconstruction gaps, and zero- and shuffle-perturbation interventions --
we demonstrated that all sampled-latent variants trained with
dense-view decoder conditioning suffer a complete posterior collapse:
the KL vanishes, the prior forgets the context, and the decoder becomes
independent of the latent. We traced this failure to a train/deploy
information mismatch, i.e., privileged dense views reaching the decoder
directly during training, and derived three orthogonal, actionable
remedies: context-conditioned decoding (now the default in the released
configuration), free-bits floors, and deployment-guided
prior-reconstruction supervision.

\textbf{Future work} includes (i) validating the collapsed-vs-healthy
regime under the context-conditioned decoder at the full 400-epoch
protocol, (ii) extending P5 to hierarchical local latents for
fine-grained uncertainty maps, (iii) calibrating the predictive
variance against real reconstruction errors for use in dose-guided
adaptive radiotherapy, and (iv) transferring the diagnostic battery to
other modality-paired inverse problems (e.g., sparse-view MRI,
low-dose PET) where the information-flow argument applies verbatim.

\section*{Acknowledgements}

The authors gratefully acknowledge the computing resources and the
LUNA16 challenge data \cite{setio2017luna16}. The authors declare no
conflict of interest.